%==============================================================
% BAB I PENDAHULUAN
%==============================================================

\section{Pendahuluan}

\subsection{Latar Belakang dan Urgensi Masalah}

[Tuliskan latar belakang permasalahan di sini. Jelaskan konteks kesehatan di Indonesia, beban ganda penyakit tropis dan penyakit tidak menular, serta hambatan akses di fasilitas kesehatan tingkat pertama (FKTP) dan wilayah kepulauan/pedalaman. Cantumkan data pendukung dan referensi seperti \citep{Kemenkes2024}].

\begin{enumerate}
    \item \textbf{[Masalah 1 - Antrean dan Administrasi FKTP]:} [Uraikan kendala antrean panjang di Puskesmas, kasus ringan SKDI 4A yang seharusnya bisa mandiri, dan pembatalan tes lab karena pasien tidak puasa].
    \item \textbf{[Masalah 2 - Akses Geografis dan Salah Rujukan]:} [Uraikan kendala transportasi wilayah kepulauan/pedalaman dan risiko salah rujukan ke faskes tanpa stok obat program].
    \item \textbf{[Masalah 3 - Rendahnya Skrining Mandiri Usia Produktif]:} [Uraikan rendahnya pemeriksaan berkala tensi, gula darah, dan antropometri di masyarakat tanpa alat medis].
    \item \textbf{[Masalah 4 - Bahaya Swamedikasi dan Tanggap Darurat]:} [Uraikan risiko konsumsi obat bebas keliru seperti NSAID pada demam dengue dan keterlambatan akses darurat].
\end{enumerate}

\subsection{Rumusan Masalah}

[Rumuskan pertanyaan kunci yang dijawab oleh inovasi ini]:
\begin{enumerate}
    \item Bagaimana mengembangkan sistem triase mandiri berbasis *Medical Large Language Model* yang terikat bukti klinis (*evidence-grounded*) untuk memilah keparahan kondisi tanpa risiko halusinasi?
    \item Bagaimana merancang integrasi direktori fasilitas kesehatan nasional SATUSEHAT MSI dan rujukan terdekat berbasis privasi tanpa meminta data kependudukan sensitif?
    \item Bagaimana menyediakan deteksi interaksi obat aktif dan proksi antropometri fisik bagi masyarakat yang tidak memiliki instrumen medis rumah tangga?
\end{enumerate}

\subsection{Tujuan Inovasi}

[Tuliskan tujuan spesifik dan terukur dari inovasi yang diusulkan]:
\begin{enumerate}
    \item Mewujudkan sistem triase mandiri pra-hospital 4 rute (Merah, Kuning Khusus, Kuning, Hijau) berstandar PPK Dokter di FKTP Kemenkes RI.
    \item Mengotomasi navigasi rujukan faskes terdekat ber-WhatsApp dan transmisi ringkasan klinis terenkripsi (*QR Code*).
    \item Memberikan edukasi keamanan swamedikasi obat dan kepatuhan persiapan pra-laboratorium bagi pasien.
\end{enumerate}
